% Compile this document in Overleaf using XeLaTeX.
% Menu -> Compiler -> XeLaTeX -> Recompile
\documentclass[12pt,a4paper]{article}
\usepackage[margin=1in]{geometry}
\usepackage{fontspec}
\usepackage{array}
\setmainfont{Latin Modern Roman}
\newfontfamily\devanagarifont[Script=Devanagari]{Noto Serif Devanagari}
\usepackage{hyperref}
\title{Typesetting in Indian Languages\\\large Practical 8 and Practical 9}
\author{LaTeX Practical Study Notes}
\date{ }

\begin{document}
\maketitle
\tableofcontents
\newpage

\section{Basic terminology}
\begin{description}
  \item[Language] The language being written, such as Hindi, Sanskrit, Punjabi or Malayalam.
  \item[Script] The writing system used, such as Devanagari, Gurmukhi or Malayalam script.
  \item[Font] The typeface used to draw characters. Rachana, for example, is a Malayalam font.
  \item[Input scheme] The method used to enter characters, such as Unicode, Velthuis or ITRANS.
\end{description}

\section{Practical 8: Devanagari using Devnag / Velthuis}
This practical uses a transliteration source file and the legacy \texttt{devnag} preprocessor.

\subsection{Procedure}
\begin{enumerate}
  \item Write the transliteration source in a file with the \texttt{.dn} extension.
  \item Mark Devanagari text using the command \verb|{\dn ...}| and load the Devanagari package in the appropriate LaTeX source.
  \item Run the Devnag preprocessor to generate a \texttt{.tex} file.
  \item Compile the generated \texttt{.tex} file using the required LaTeX engine and font resources.
\end{enumerate}

\subsection{Example}
The transliteration example is:
\begin{verbatim}
{\dn ga.nita}
\end{verbatim}
It represents the Devanagari word:

\begin{center}
  {\devanagarifont गणित}
\end{center}

\noindent\textbf{Important:} The Devnag workflow requires the legacy Devnag program and its font resources. The command \verb|{\dn ...}| is input for that workflow; it is not a general-purpose command available in every Overleaf project.

\section{Practical 9: Indian-language typesetting using ITRANS}
ITRANS is a separate transliteration preprocessor. Its directives belong to ITRANS input, not ordinary LaTeX.

\subsection{Example ITRANS input}
Create a source file such as \texttt{input.itx}:
\begin{verbatim}
#indianifm=dvng.ifm
#indianfont=\fransdvng
#indian
... transliterated words ...
#endindian
\end{verbatim}

\subsection{Procedure}
\begin{enumerate}
  \item Write the transliterated words and ITRANS directives in a \texttt{.itx} file.
  \item Run the ITRANS preprocessor:
  \begin{verbatim}
itrans -i input.itx -o output.tex
  \end{verbatim}
  \item Compile the generated file:
  \begin{verbatim}
pdflatex output.tex
  \end{verbatim}
\end{enumerate}

\section{Devnag and ITRANS: difference}
\begin{center}
\begin{tabular}{|>{\raggedright\arraybackslash}p{0.22\linewidth}|p{0.32\linewidth}|p{0.32\linewidth}|}
\hline
 & \textbf{Devnag / Velthuis} & \textbf{ITRANS} \\
\hline
Input file & Usually \texttt{.dn} & Usually \texttt{.itx} \\
\hline
Preprocessor & Devnag & ITRANS \\
\hline
Output & Generated LaTeX source & Generated LaTeX source \\
\hline
Key point & Uses Devnag transliteration conventions & Uses ITRANS conventions and directives \\
\hline
\end{tabular}
\end{center}
The two preprocessors are different. Their transliteration conventions and commands are not interchangeable.

\section{Modern Unicode alternative}
Modern XeLaTeX or LuaLaTeX can typeset Unicode Indic text when a suitable font is available. This example uses XeLaTeX and the Noto Serif Devanagari font:
\begin{verbatim}
\usepackage{fontspec}
\newfontfamily\devanagarifont[Script=Devanagari]{Noto Serif Devanagari}
{\devanagarifont गणित}
\end{verbatim}

\noindent This is a modern Unicode workflow. It is useful for reading and typing Indic text in Overleaf, but it is not a replacement for demonstrating the syllabus's legacy Devnag and ITRANS preprocessing steps if those are specifically required in the practical.

\section{Quick viva revision}
\begin{enumerate}
  \item \textbf{What is a script?} A writing system used to represent a language.
  \item \textbf{What is transliteration?} Representing text from one script using the characters or conventions of another.
  \item \textbf{What does Devnag do?} It preprocesses Devnag/Velthuis transliteration input into LaTeX source.
  \item \textbf{What does ITRANS do?} It converts ITRANS transliteration input into output suitable for LaTeX processing.
  \item \textbf{Can Devnag and ITRANS be used interchangeably?} No. They are separate preprocessors with different conventions.
  \item \textbf{Which engine supports Unicode Indic text with fontspec?} XeLaTeX or LuaLaTeX.
\end{enumerate}

\end{document}
